%% notes-tex/publishing-system/sections/3-notes-tex.tex

\section{A typeset note\secstatus{todo}}
\label{sec:notes}

%% SOURCE: notes/publishing-system.mermaid.notes-tex-build.md, rendered by
%% notes/assets/publish/scripts/mermaid_pdf.sh and scaled to the text block by
%% `make diagrams`; the targets and outputs are those of notes-tex/common/Makefile.common.
\tcfig{figures/notes-tex-build.pdf}{What \texttt{make} builds for a typeset note.
Same colors as Figure~\ref{fig:paper}; gray is the gate, which builds nothing.
Figures and tables enter from the experiment scripts on the left; the prose is
the only thing edited in the document directory.}{fig:notes}

\notesec{Layout of a document directory}
\texttt{<slug>.tex} is the root and is named for the directory, so the
\texttt{eqtl-data-model} directory builds \texttt{eqtl-data-model.pdf}; four
documents named \texttt{main.pdf} would be four files a file-opener cannot tell
apart. \texttt{sections/*.tex} hold the prose, one file per top-level section.
\texttt{tables/*.tex} are written by scripts and never edited by hand.
\texttt{figures/*.pdf} are converted from script-written SVG panels or exported
from draw.io. \texttt{references.bib} is pulled, not written. The Makefile sets at
most a collection key and the image directories to read, then includes the common
file. \texttt{tcdoc.sty} is symlinked into the directory at build time, because
Tectonic does not honor \texttt{TEXINPUTS}; the link is ignored by git and any
build recreates it.

%% SOURCE: hand-authored from notes-tex/common/Makefile.common and the
%% 025-additive-baselines Makefile.
\begin{table}[htbp]
\centering
\footnotesize
\caption{Targets every typeset note has, from the common Makefile, plus the two a
document may add.}
\label{tab:notetargets}
\begin{tabular}{lp{0.66\textwidth}}
\toprule
Target & Produces \\
\midrule
\texttt{make} & \texttt{<slug>.pdf}, the draft view with status chips and source pointers \\
\texttt{make clean-view} & \texttt{<slug>-clean.pdf}, chips and pointers hidden, provenance flags kept \\
\texttt{make check} & nothing; the formatting, provenance, citation and style gate on the last build \\
\texttt{make plots} & \texttt{figures/{}NAME.pdf} from each referenced \texttt{notes/assets/images/<slug>/NAME.svg} \\
\texttt{make figures} & \texttt{figures/{}NAME.pdf} from each referenced draw.io source \\
\texttt{make bib-pull} & \texttt{references.bib} for this document from the served store \\
\texttt{make bib}, \texttt{bibcheck} & the same through Better BibTeX on a machine running Zotero; the collision report \\
\texttt{make watch} & rebuild on save \\
\midrule
\texttt{make short} & \texttt{<slug>-short.pdf}, a second root over the same sections (025-additive-baselines) \\
\texttt{make diagrams} & mermaid figures rendered and scaled to the text block (this document) \\
\bottomrule
\end{tabular}
\end{table}

\notesec{The two views}
The draft view carries the status chips in every heading, and so in the contents,
which is the status board: \texttt{todo} may be edited freely, \texttt{tent} has
been read by the author and is kept stable, \texttt{final} is ready. It also
prints a gray source pointer under generated content. \texttt{make clean-view}
rebuilds with the \texttt{clean} package option, which removes both, and keeps the
provenance flags and their key, because an unverified number must not lose its
flag when the document leaves the group.

\notesec{Converting a panel}
A plot script writes its panel as a true-size SVG whose width is in draw.io's unit
of 100 per inch, a number no renderer interprets the same way twice: librsvg has
read it as points in one version and as CSS pixels in another, so one 179 mm panel
converted to a 249 mm PDF on one machine and a 186 mm PDF on another. \texttt{make
plots} therefore measures rather than assumes: it converts a 100-unit probe square,
reads the size that came out, and converts every panel at the zoom that maps the
probe to 72 points. The result is the panel at its true millimeters, placed with
\texttt{\textbackslash tcfig} at 1:1 so 6 pt type in the SVG is 6 pt on the page.

\notesec{The gate}
\texttt{make check} reads the last build's log and every \texttt{.tex} file and
reports six things: any heading without a status chip; any referenced figure wider
than the 182 mm text block; any \texttt{\textbackslash cite} key absent from
\texttt{references.bib}; any file with a table but no \texttt{\%\% SOURCE:} line;
British spelling and em-dashes in prose, with narration and self-reference as
warnings; and overfull boxes above 5 pt from the log, which is how a table running
off the page is caught without anyone reading page 8. Errors exit non-zero,
warnings do not, since a \texttt{todo} section is allowed to be incomplete.
Trailing restatement is not checked; an attempt ran at one true positive in ten and
was removed.

\notesec{A second build from the same sections}
A document that needs a short version for sharing adds a second root that sets a
switch, \texttt{\textbackslash longfalse}, and wraps the long-only material in
\texttt{\textbackslash iflong \ldots \textbackslash fi}. Its Makefile adds
\texttt{make short} and \texttt{make check-short}; nothing in the common file
changes, and the short PDF publishes with \texttt{--pdf <slug>-short}.
